\documentclass[10pt,a4paper]{amsart}
\usepackage[margin=19mm]{geometry}
\usepackage{amsmath,amssymb,mathtools}
\usepackage[scr=rsfs,cal=euler]{mathalfa}
\usepackage{xfrac}
\usepackage[unicode,hidelinks,bookmarks=false]{hyperref}
\newcommand{\ie}{i.e.\ }
\newcommand{\cf}{cf.\ }
\newcommand{\resp}{respectively\ }
\newcommand{\Subsection}{\subsection}
\providecommand{\nobreakdash}{\nobreak\hbox{-}}
\setlength{\emergencystretch}{2em}
\multlinegap0pt
\usepackage{bm}
\usepackage{bbm}
\usepackage{dsfont}
\usepackage{xspace}
\usepackage{cancel}
\usepackage{mathtools}
\usepackage[shortlabels]{enumitem}
\usepackage{amsthm}
\makeatletter
\@ifundefined{theorem}{\newtheorem{theorem}{Theorem}}{}
\makeatother
\makeatletter
\@ifundefined{proof}{\renewenvironment{proof}{{\noindent\bfseries Proof.}}{\qed}}{}
\makeatother
\title[Explicit Formulas and Unimodality Phenomena for General Posi]{Explicit Formulas and Unimodality Phenomena for General Position Polynomials}
\author{Bilal Ahmad Rather}
\date{Source: arXiv:2603.06930v3 (CC BY 4.0)}
\begin{document}
\maketitle

\noindent\emph{Source and licence.} Explicit Formulas and Unimodality Phenomena for General Position Polynomials. Version arXiv:2603.06930v3 (CC BY 4.0), distributed under the Creative Commons Attribution 4.0 International licence, as recorded in the arXiv licence metadata for this exact version and shown on the abstract page https://arxiv.org/abs/2603.06930v3. Full rights notice: licenses/arxiv-2603.06930-CC-BY-4.0.txt.\par\smallskip

\noindent\textit{Editorial scope.} Counted unit: the Theorem whose source label is \texttt{thm:Tstar-shape} in the article (source0.tex lines 1023--1038, source label \texttt{thm:Tstar-shape}), delivered together with the article's own complete original proof of it (source0.tex lines 1040--1093, one \texttt{proof} environment of 54 lines). No mathematics is written, repaired or re-derived: statement and proof are the article's own text line for line. Required context retained uncounted: thm:correct-gpa-Tstar (lines 843--896). Every definition, display and label of those lines is retained unchanged. Omitted from this package and not redistributed: 1--842, 897--1022, 1039--1039, 1094--1473. Nothing inside a retained excerpt is omitted or altered: every retained body is delivered exactly as the article's lines are written, so no omission receipt of any kind accompanies this package. Citations: the retained excerpts carry no citation command at all, so no bibliography entry of the article is transcribed and no reference is redistributed. Reference adaptation: none was needed and none was made; every reference inside the delivered unit resolves inside it, so no unresolved reference or citation is produced. Printed slips kept literal: no printed word, symbol, hypothesis or number of the counted statement or of its proof was altered. Layout and class: the article's own class is \texttt{article} with packages including graphicx, graphics, amsthm, amsfonts, amsbsy, amssymb, amsmath, cite, array, arydshln, hyperref, float, tikz, inputenc; the wrapper is the lane's pinned \texttt{amsart} editorial wrapper at 10pt on A4, which restates the theorem-like environments the retained text needs and adds the article's own macro definitions that the retained text needs (none), with extra wrapper packages (bm, bbm, dsfont, xspace, cancel, mathtools, enumitem). The delivered numbering is the unit's own. Assets: no retained span contains a figure environment, a graphics-inclusion command, a PDF-page inclusion, a table or a TikZ picture; no image file is copied into this package, the package declares no asset dependency and no image file is redistributed with it. Licence: version arXiv:2603.06930v3 of this article is distributed under the Creative Commons Attribution 4.0 International licence, as recorded in the arXiv licence metadata for this exact version and shown on the abstract page https://arxiv.org/abs/2603.06930v3. A full rights notice is delivered at licenses/arxiv-2603.06930-CC-BY-4.0.txt. Characters: every retained excerpt is pure ASCII, so the delivered engine, XeLaTeX with its default Latin Modern text font, reports no missing character.
		\begin{theorem}\label{thm:correct-gpa-Tstar}
			Let $T^*(r,a)$ be the graph with vertex set
			$
			V\bigl(T^*(r,a)\bigr)=\{ i_j:\ i\in [a],\ j\in [r] \},
			$
			where two distinct vertices $i_j$ and $k_\ell$ are adjacent if and only if
			$
			i\neq k $ and $ j\neq \ell.
			$
		If $\psi\bigl(T^*(r,a)\bigr)=\sum_{k\ge 0}\alpha_k x^k,$ then $\alpha_0=1$, $\alpha_1=ar$, and
			$
			\alpha_2=\binom{ar}{2}.
			$
			Moreover, the  GPA of $T^*(r,a)$ is given by the following three cases.
			\begin{enumerate}[label=\textup{(\roman*)}]
				\item If $\min\{r,a\}=1$, then $T^*(r,a)\cong \overline{K}_{ar}$ and hence
			$\psi\bigl(T^*(r,a)\bigr)=(1+x)^{ar}. $
			\item If $\min\{r,a\}=2$, and $m=\max\{r,a\}$, then $T^*(r,a)\cong K_{m,m}-M$, where $M$ is a perfect matching, and
			$$
			\psi\bigl(T^*(r,a)\bigr)
			=
			1+2mx+\binom{2m}{2}x^2+2\sum_{k=3}^{m}\binom{m}{k}x^k.
			$$
			\item If $r\ge 3$ and $a\ge 3$, then for every $k\ge 3$,
			$$
			\alpha_k
			=
			a\binom{r}{k}
			+
			r\binom{a}{k}
			+
			k!\binom{a}{k}\binom{r}{k}
			+
			\binom{a}{2}\binom{r}{2}\binom{4}{k},
			$$
			and
			{\footnotesize$$
			\psi\bigl(T^*(r,a)\bigr)
			=
			1+arx+\binom{ar}{2}x^2
			+
			\sum_{k=3}^{\max\{a,r,4\}}
			\left(
			a\binom{r}{k}
			+
			r\binom{a}{k}
			+
			k!\binom{a}{k}\binom{r}{k}
			+
			\binom{a}{2}\binom{r}{2}\binom{4}{k}
			\right)x^k.
			$$}
			\end{enumerate}
		\end{theorem}

		\begin{theorem}\label{thm:Tstar-shape}
			Let $\psi\bigl(T^*(r,a)\bigr)=\sum_{k\ge 0}\alpha_kx^k$	be the  GPA of $T^*(r,a)$ as in Theorem~\ref{thm:correct-gpa-Tstar}. Then the following hold. 
			\begin{enumerate}[label=\textup{(\roman*)}]
				\item If $\min\{r,a\}=1$, then $\psi\bigl(T^*(r,a)\bigr)=(1+x)^{ar},$ hence the coefficient sequence is real-rooted, log-concave, and unimodal.
				
				\item If $\min\{r,a\}=2$, and $m=\max\{r,a\}$, then
				$$
				\psi\bigl(T^*(r,a)\bigr)
				=
				1+2mx+\binom{2m}{2}x^2+2\sum_{k=3}^{m}\binom{m}{k}x^k.
				$$
				In this case the coefficient sequence is: log-concave if and only if $m\le 4$,  unimodal for all $m\neq 8$, and   not unimodal for $m=8$.
				
				\item If $r\ge 3$ and $a\ge 3$, then there is no uniform behavior in general, as the family $T^*(r,a)$ contains examples whose  GPAs are log-concave and unimodal,  unimodal but not log-concave, and  neither unimodal nor log-concave. In particular, all three phenomena already occur in the subfamily $T^*(r,3)$.
			\end{enumerate}
		\end{theorem}

		\begin{proof}
			For $\min\{r,a\}=1$, the graph $T^*(r,a)$ is edgeless on $ar$ vertices, so $\psi\bigl(T^*(r,a)\bigr)=(1+x)^{ar}.$
			Hence, its coefficients are binomial coefficients, and therefore log-concave and unimodal. Now, assume $\min\{r,a\}=2$, and write $m=\max\{r,a\}$. So, by Theorem~\ref{thm:correct-gpa-Tstar}, we have
			$\alpha_0=1, \alpha_1=2m, \alpha_2=\binom{2m}{2}=m(2m-1),$	and $\alpha_k=2\binom{m}{k}\qquad (k\ge 3).$  For $k\ge 4$, the tail $(\alpha_k)_{k\ge 3}$ is a positive scalar multiple of the binomial sequence, hence it is log-concave. Therefore, only the indices $k=1,2,3$ need to be checked. A direct computation gives
			$$
			\alpha_1^2-\alpha_0\alpha_2
			=
			(2m)^2-\binom{2m}{2}
			=
			m(2m+1)>0,
			$$
			$$
			\alpha_2^2-\alpha_1\alpha_3
			=
			\bigl(m(2m-1)\bigr)^2-(2m)\cdot 2\binom{m}{3}
			=
			\frac{m^2}{3}\bigl(10m^2-6m-1\bigr)>0,
			$$
			and
			$$
			\alpha_3^2-\alpha_2\alpha_4
			=
			\left(2\binom{m}{3}\right)^2-m(2m-1)\cdot 2\binom{m}{4}
			=
			-\frac{m^2(m-1)(m-2)}{36}\bigl(2m^2-9m+1\bigr).
			$$
			Hence $\alpha_3^2\ge \alpha_2\alpha_4$ if and only if $2m^2-9m+1\le 0$, that is, if and only if $m\le 4$.  For unimodality, we note that 
			$$
			\alpha_2-\alpha_3
			=
			m(2m-1)-2\binom{m}{3}
			=
			-\frac{m}{3}\bigl(m^2-9m+5\bigr).
			$$
			Thus, $\alpha_2\ge \alpha_3$ for $m\le 8$, while $\alpha_2<\alpha_3$ for $m\ge 9$. If $2\le m\le 7$, then the binomial tail $2\binom{m}{k}$ is nonincreasing for $k\ge 3$, since the binomial coefficients of order $m$ reach their maximum at $k\le 3$. Since also $\alpha_2\ge \alpha_3$, we get $	\alpha_0<\alpha_1<\alpha_2\ge \alpha_3\ge \alpha_4\ge \cdots,
			$ and the sequence is unimodal. If $m=8$, then
			$$
			\psi\bigl(T^*(8,2)\bigr)=1+16x+120x^2+112x^3+140x^4+112x^5+56x^6+16x^7+2x^8,
			$$
			whose coefficients satisfy $120>112<140,$ so the sequence is not unimodal. If $m\ge 9$, then $\alpha_2<\alpha_3$, and the tail $(2\binom{m}{k})_{k\ge 3}$ is itself unimodal. Since, $\alpha_1<\alpha_2<\alpha_3$, the whole coefficient sequence is unimodal.  Finally, for $r\ge 3$ and $a\ge 3$,  Theorem~\ref{thm:correct-gpa-Tstar} shows that the coefficient sequence is a sum of four different contributions 
			$$
			\alpha_k
			=
			a\binom{r}{k}
			+
			r\binom{a}{k}
			+
			k!\binom{a}{k}\binom{r}{k}
			+
			\binom{a}{2}\binom{r}{2}\binom{4}{k}
			\qquad (k\ge 3).
			$$
			These, contributions interact in different ways for different pairs $(r,a)$, and this leads to all three types of behavior listed in part~(iii). 
		\end{proof}

\end{document}
